\section{Conclusion and Future Work}
\label{sec:conclusion}

\subsection{Conclusions}

In this work, we investigated a cut-and-paste outlier exposure strategy for improving out-of-distribution (OOD) detection in semantic segmentation models. The proposed approach leverages real object instances from COCO and integrates them into road-scene images using structured augmentation priors, namely scale, spatial, and categorical priors.

The fine-tuning procedure combines standard semantic segmentation supervision with an additional OOD objective applied to synthetic anomaly regions. This formulation allows the model to be explicitly exposed to object-level anomalies while preserving its ability to perform in-distribution segmentation, aided by the use of a frozen encoder.

Experimental results across multiple benchmarks, including Lost\&Found, Fishyscapes, Road Anomaly, Road Anomaly 21, and Road Obstacle 21, show changes in OOD detection performance after fine-tuning. Improvements are more consistent for medium and large-scale anomalies, while results for small-scale anomalies remain limited in most cases. Performance varies across datasets, indicating that the effectiveness of the approach depends on scene structure and anomaly characteristics.

Overall, the study suggests that cut-and-paste outlier exposure can be a practical mechanism for improving OOD awareness in segmentation models without architectural modifications, although its impact is not uniform across all evaluation settings.

\subsection{Future Work}

Several directions could be explored to extend this work.

First, the current cut-and-paste mechanism could be improved by incorporating more realistic compositing techniques. In particular, learning-based blending or diffusion-based inpainting could reduce the domain gap between synthetic and real anomalies.

Second, the categorical prior is currently defined using either fixed weights or simple adaptive heuristics. More principled approaches could be explored, such as learning the sampling distribution jointly with the model or conditioning it on scene context.

Third, the spatial prior relies on handcrafted perspective assumptions. This could be replaced or complemented by depth estimation models to provide more accurate geometric placement of objects.

Fourth, while the encoder is frozen in this work to preserve pretrained representations, partial or progressive unfreezing strategies could be investigated to assess potential gains in adaptation capacity.

Finally, evaluation could be extended beyond current benchmarks to include more diverse and real-world driving scenarios, particularly those with complex weather conditions or long-tail object distributions, to better assess robustness in practical deployments.